\chapter*{Abstract}
\addcontentsline{toc}{chapter}{Abstract}

The Internet forwards packets without regard for why they were sent, so the network optimizes a target
that rarely matches what an application actually values. Path-aware architectures such as SCION change
this by letting endpoints choose among multiple end-to-end paths, placing the path decision in the hands
of the one party that knows the application's intent. This thesis studies how to make that choice well:
how to select the best path -- or the best combination of paths -- given both the application's intent and
the current state of the network. We build two learned controllers to answer it, at two granularities:
a single-path selector for general flows, and a joint controller for real-time video that decides how
much traffic each path carries \emph{and} what the application produces. Both take a single idea from
recent congestion-control work -- that application intent should be a first-class input which reshapes
a network mechanism rather than a constant compiled into it -- and carry it one layer deeper than it
has been taken before.

The strongest result is that the two decisions cannot be made well apart. A two-timescale pair of
reinforcement-learning agents co-optimizing a video call's encoding bitrate and its per-path split over
Multipath QUIC beats all eight comparison points under all three application intents -- including
minRTT, the scheduler every Multipath QUIC implementation ships, and a WebRTC-style deployed stack --
by 0.065 to 0.143 QoE over the strongest of them, three to seven times the spread between training
runs of the same method, and removing either agent forfeits most of the achievable quality. What makes
this possible is architectural rather than incidental: a permutation-equivariant head that scores each
path with shared weights, so one policy serves any number of paths. Substituting a conventional
fixed-dimension head costs 0.231 QoE and drops the system from first place in nine of nine evaluation
cells to none, which makes it the largest measured effect in this thesis, and the substitute cannot even
be loaded at a path count other than the one it trained on, where the shared-weight policy runs
unmodified from two paths to twelve at a parameter count that does not move. The same
architecture carries the single-path study, where it raises residual bottleneck capacity by 41\% over a
fixed-action control at less than half the parameter count, and brings a learned selector to within
0.002 reward of a per-context oracle while probing one path per decision instead of thirty.

Intent conditioning is the thread that connects the two, and its honest form is a design lemma with a
measured boundary. We show analytically that \emph{where} the conditioning enters decides whether it
has any effect at all: a signal reaching only the terms shared by every candidate cancels under the
comparison that selects a path, exactly and on the entire action when that comparison is a softmax,
however much the policy is trained on it. Routing it into the per-candidate comparison instead lets one
policy serve a range of intents without retraining, including intents it never trained on, and in the
multipath setting lets one policy replace three separately trained specialists at a third of the
training cost with no penalty detectable at the resolution our experiments achieve. But the guarantee
is narrower than the behavior. The cancellation is exact only for scorers additively separable in the
conditioning, and we measure three ways an ordinary deployed architecture violates that premise, so
that a trained nonlinear encoder routes intent past a cancellation the theory says should stop it. The
principle survives as a statement about which routings carry a guarantee, not which happen to work.
Both settings also bound how much such conditioning can be worth: intents must genuinely conflict, and
in the single-path environment loss is near zero on 99\% of selections, leaving goodput against latency
as the one axis on which alignment can be demonstrated.

We report where the results do not reach, and two of those findings are more useful than what they
replaced. What governs the payoff of joint control is not whether the network is moving but how
unevenly capacity is spread across the paths and how strictly the application prices delay -- a
stationary but asymmetric network defeats application-only adaptation as thoroughly as a churning one,
and the two axes interact rather than adding, the three intents differing fourteen times more on a harsh
network than on a benign one. The same map bounds the headline claim: it holds on a milder held-out
network and fails on a harsher one, where a hand-written scheduler wins every column.
And on a packet-level backend the ranking transfers for three of six independently trained policies and
not the other three, with analytical-backend score uncorrelated with which, so a cheap surrogate can be
an adequate training environment while being useless for selecting among the policies it produces. Four
results in this thesis were reported from a single training run and did not survive being run three
times; we report that pattern as a result, because on this problem the spread between training runs
exceeds the confidence interval over evaluation episodes by an order of magnitude. Together the two
parts argue that on a path-aware Internet, application intent should steer every controllable decision
between the application and the wire, and they chart what remains to make such control deployable.

\gap{one scope sentence to add once you are happy with the wording above: the multipath study assumes
the path multiplicity a path-aware architecture provides rather than running on a SCION data plane --
there is no beaconing, path store, or segment composition in \cref{ch:multipath}, and the six-path
topology is hand-specified (\cref{tab:app:p2topo}). Given the title, a SCION-group reader will ask, and
it is better answered here in one clause than discovered in \cref{sec:mp:impl}.}
